% Generated draft: requires mathematical and source review.
\documentclass[11pt]{article}
\usepackage[margin=25mm]{geometry}
\usepackage{fontspec}
\setmainfont{Latin Modern Roman}
\usepackage{amsmath,amssymb,mathtools,unicode-math}
\setmathfont{Latin Modern Math}
\usepackage{enumitem,needspace,xurl,hyperref}
\setlength{\parindent}{0pt}
\setlength{\parskip}{5pt}
\setlength{\emergencystretch}{3em}
\newcommand{\N}{\mathbb{N}}
\newcommand{\Z}{\mathbb{Z}}
\newcommand{\Q}{\mathbb{Q}}
\newcommand{\R}{\mathbb{R}}
\newcommand{\C}{\mathbb{C}}
\newcommand{\F}{\mathbb{F}}
\newcommand{\A}{\mathbb{A}}
\newcommand{\PP}{\mathbb{P}}
\newcommand{\E}{\mathbb{E}}
\newcommand{\eps}{\varepsilon}
\newcommand{\dd}{\,\mathrm{d}}
\begin{document}
\section*{UnsolvedMath ID: 30005562}
\textbf{Source problem code:} catalogue entry
\subsection{Definitions and mathematical statement}
For a field \(K\), let \(K^{\mathrm{sep}}\) be a separable closure and let \(G_K=\operatorname{Gal}(K^{\mathrm{sep}}/K)\) be its absolute Galois group. Using the usual Galois-cohomological convention, define
\[
\operatorname{cd}(K)=\inf\bigl\{d\in\Z_{\geq0}:H^q_{\mathrm{cont}}(G_K,A)=0\text{ for every discrete torsion }G_K\text{-module }A\text{ and every integer }q>d\bigr\},
\]
with value \(\infty\) if the set is empty.

In the usual neighborhood-based terminology, a Stein compact set is a compact subset \(S\) of a complex analytic space \(X\) that has a basis of Stein open neighborhoods in \(X\). A possible meaning of meromorphic functions along \(S\) is the ring of meromorphic germs
\[
\mathcal M(S)=\varinjlim_{U\supseteq S}\Gamma(U,\mathcal M_X),
\]
where \(U\) ranges over open neighborhoods and \(\mathcal M_X\) is the sheaf of meromorphic functions. The supplied record does not specify this construction, the ambient space, or hypotheses ensuring that this ring is a field. It also does not specify how the complex dimension of the compact subset is defined. These conventions therefore require confirmation from the source.

\textbf{Intended question, subject to unresolved hypotheses.} For every Stein compact set \(S\) of complex dimension \(n\), with \(\mathcal M(S)\) understood as the meromorphic function field intended in the source, is it true that
\[
\operatorname{cd}\bigl(\mathcal M(S)\bigr)=n?
\]
Here \(n\) is the nonnegative integer described as the complex dimension in the supplied statement. No restriction excluding dimension zero is supplied. The record provides no connectedness, normality, irreducibility, reducedness, purity-of-dimension, or ambient hypotheses for this question. Consequently, the displayed question records the intended target, but does not determine a fully specified universal assertion until the meaning of dimension and the field construction are resolved. The question form preserves the original excerpt; the reconstructed statement's instruction to prove the equality is not evidence that the equality is established.
\subsection{Short English statement}
The question asks whether the Galois cohomological dimension of the meromorphic function field associated with a Stein compact set equals its complex dimension. The supplied record leaves essential geometric hypotheses and the field construction unspecified, so this is a statement-only draft requiring mathematical review.
\subsection{Sources}
\begin{itemize}
\item Supplied local UnsolvedMath registry record, permanent ID \(30005562\), title \emph{Cohomological Dimension of Stein Meromorphic Function Fields}; exact source problem code \texttt{catalogue entry}. The record belongs to the Oberwolfach Reports Open Problems collection and gives proposed year 2023. Its original excerpt associates the question with Olivier Benoist's contribution \emph{On the field of meromorphic functions on a Stein compact set}; no separate proposer is recorded.
\item \emph{Recent Trends in Algebraic Geometry} (2023), Oberwolfach Reports, pp.~1485--1542, DOI \texttt{10.4171/owr/2023/27}: \url{https://doi.org/10.4171/owr/2023/27}. This is the supplied report provenance. The registry labels the reconstructed statement \texttt{corrected\_verified} and says it was checked against the official report; that is a supplied verification claim, not an independent verification in this draft.
\item Benoist, \emph{On the field of meromorphic functions on a Stein surface}: \url{https://doi.org/10.1515/crelle-2025-0071}. This link is supplied in the literature sources.
\item Benoist, \emph{On the field of meromorphic functions on a Stein compact set}: \url{https://doi.org/10.4171/owr/2023/27}. This supplied literature link identifies the workshop report rather than a separately pinpointed contribution.
\item Literature assessment dated exactly \texttt{2026-08-22}: supplied dated curation, not independently verified. The record labels the problem partially solved and reports that the equality is proved for connected normal Stein curves and, by what it calls Benoist's 2026 theorem, for Stein surfaces, while the general question in dimension \(n\geq3\) remains open. The background explicitly describes this as dated literature triage derived from the supplied OWR curation and requests independent verification of time-sensitive status and cited claims. These claims are retained as source qualifications, not asserted here as verified mathematical facts.
\end{itemize}
\end{document}
